\section{Blow-ups and the normal-meridian criterion}
\label{sec:blowups}

Call a normal-crossing pair $(X,D)$ \emph{good} when every connected
stratum, including $U=X\setminus D$, is aspherical and every full
normal-meridian lattice $\Lambda_S\to\pi _1(U)$ is injective. Branches
are taken on the normalization; a self-intersection contributes two
local branches. A \emph{stratal centre} is a smooth closed centre
which, in every normal-crossing chart, is the intersection of a
subset of the boundary branches. Disjoint unions of such centres
are allowed. This local condition is more precise than demanding
that the centre be one globally irreducible stratum closure.

\begin{theorem}[The exact blow-up criterion; \status{Proved}]
\label{P2a}
Let $(X,D)$ be good, and let $Z\subset D$ be a smooth centre having
normal crossings with $D$. Write $\beta:\widetilde X=\operatorname{Bl}_Z X
\to X$ and $\widetilde D=(\beta^{-1}D)_{\rm red}$. Then
$(\widetilde X,\widetilde D)$ is good if and only if $Z$ is stratal.
A nonstratal blow-up creates an explicit integral kernel in a full
normal-meridian lattice, even when the individual new meridian is
infinite. The same kernel obstructs profinite injectivity.
\end{theorem}
\begin{proof}
At a general point of a connected component of the centre choose
coordinates in which
\[
 D=\{x_1\cdots x_s=0\},\qquad
 Z=\{x_1=\cdots=x_s=y_1=\cdots=y_t=0\}.
\]
The normal-crossing hypothesis on the centre permits precisely this
form. The centre is stratal exactly when $t=0$. Denote the old
meridians by $m_1,\ldots,m_s$. The exceptional meridian is
$m_E=m_1\cdots m_s$: a transverse exceptional circle winds once
around every $x_i$, and zero times around coordinates not defining
the boundary.

If $t>0$, use the blow-up chart with exceptional coordinate $y_1$,
where $x_i=y_1u_i$. The exceptional component and all $s$ strict
boundary transforms meet along $y_1=u_1=\cdots=u_s=0$. Their
rank-$(s+1)$ normal lattice maps into the old rank-$s$ lattice by
\[
 (a_1,\ldots,a_s,b)\longmapsto(a_1+b,\ldots,a_s+b).
\]
Its kernel is the primitive vector $(-1,\ldots,-1,1)$, so the
composite map to $\pi _1(U)$ cannot be injective. Its image in
$\widehat\Z^{s+1}$ is still nonzero. This argument does not require
the old pair to be good, and does not assert that $m_E$ itself dies.

If $t=0$, in the chart indexed by $j$ the exceptional ray replaces
$e_j$ by $e_1+\cdots+e_s$, and the other old coordinate rays remain.
This is a unimodular change of basis. Every smaller new cone gives
a primitive sublattice of that same old lattice. Consequently all
new full normal lattices inject into $\pi _1(U)$. Over an old
stratum, each new stratum is, after a finite cover labelling the
branches, a bundle with fibre a complex torus $(\C^*)^a$. The
fibration homotopy sequence proves asphericity from that of the old
stratum. Finite covers do not alter higher homotopy groups.
Strata off the centre are unchanged. Thus the new pair is good.
\end{proof}

\begin{theorem}[Functoriality, ends, and invariance in both directions;
\status{Proved}]\label{P2b}
The blow-down induces a functor
$\Exit(\widetilde X,\widetilde D)\to\Exit(X,D)$ and is the identity
on the open stratum $U$. It preserves the pro-homotopy type of the
punctured neighbourhood of the indicated boundary. For compact
boundary and ordinary collar links, it preserves the homotopy type
of that link. For a stratal blow-up, goodness holds on one side if
and only if it holds on the other. These assertions concern the
boundary end, not unspecified additional ends of a nonproper $X$.
\end{theorem}
\begin{proof}
Each new stratum maps into an old stratum, with incidence order
preserved. Postcomposing exit paths and their simplices with
$\beta$ therefore gives the stated functor. The inverse image of
$U$ is identified with $U$, so its restriction there is the identity.

For the assertion about ends, let $V$ run through neighbourhoods
of $D$ in $X$. The punctured inverse images are literally
\[
 \beta^{-1}(V)\setminus\widetilde D=V\setminus D.
\]
They are cofinal among neighbourhoods of $\widetilde D$. Indeed,
if $W$ contains $\widetilde D$, then the proper map $\beta$ sends
the closed set $\widetilde X\setminus W$ to a closed set disjoint
from $D$; its complement is a suitable $V$. The punctured
neighbourhood pro-spaces thus agree. When a compact boundary has
a collar link, this system is represented up to homotopy by its
link, giving the additional statement.

Forward preservation was proved in Theorem \ref{P2a}. For the
reverse implication suppose the stratal blow-up is good. At a
point of an old stratum with full lattice $\Z^r$, a maximal new
cone has a unimodular rank-$r$ chart. Its injectivity consequently
proves injectivity of the old full lattice.

It remains to recover asphericity of an old stratum $S$ contained
in the centre. Pass to its finite incidence cover, if needed, to
label all $r$ normal branches and the $s$ branches defining the
centre. Let $P_S\to S$ be the punctured full normal bundle, with
fibre $(\C^*)^r$. Its homotopy sequence has connecting map
\[
 \partial:\pi _2(S)\longrightarrow\Z^r.
\]
Every element in its image dies in $\pi _1(P_S)$ and hence in
$\pi _1(U)$. The just-proved full-lattice injection forces
$\partial=0$. The open exceptional stratum $Q$ over $S$ is the
torus bundle with fibre
$(\C^*)^s/\C^*_{\rm diagonal}\cong(\C^*)^{s-1}$.
Its connecting map is the projection of $\partial$ to
$\Z^s/\Z(1,\ldots,1)$, and is therefore zero. Since $Q$ is a
stratum of the good blown-up pair, it is aspherical. The homotopy
sequence of $Q\to S$ now gives $\pi _2(S)=0$, and gives
$\pi_j(S)=0$ for $j\geq3$ because the fibre is aspherical.
This proves the assertion after descent from the incidence cover.
In particular, this argument never uses a nonexistent natural
map $Q\to U$; it uses the punctured normal bundle $P_S$.
\end{proof}

\section{Good compactifications}
\label{sec:compactifications}

\subsection{Surfaces: an exact criterion in a fixed meromorphic class}

A marked completion is a smooth compact complex surface $X$ with
reduced normal-crossing boundary $D$ and a specified isomorphism
$X\setminus D\cong U$. In the algebraic version all these data are
algebraic. Isomorphisms of marked completions induce the identity
of $U$. Two analytic marked completions belong to the same
\emph{meromorphic class} when that identity extends bimeromorphically.
This qualification is essential: a biholomorphism of noncompact
analytic surfaces need not extend meromorphically to completions.

The weighted incidence graph has one vertex for every normalized
irreducible boundary curve, decorated by its genus and
self-intersection, and one edge for every node. A self-node is a
loop, counted twice in the valency. Its \emph{core} is obtained
by repeatedly deleting genus-zero vertices of valency at most
one. Positive-genus vertices are never deleted. The complement
of the core is the \emph{fringe}; when the core is nonempty in
a connected component, its fringe consists of rational trees,
each attached to that core at one point. A \emph{corner blow-up}
is the blow-up of a boundary node, including a self-node.

\begin{lemma}[Good boundaries have no rational fringe;
\status{Proved}]\label{P3a-core-necessity}
The boundary graph of a good surface pair equals its core.
\end{lemma}
\begin{proof}
An isolated smooth rational component is a $\PP^1$ stratum and
is not aspherical. If a rational component has just one incident
branch, its open stratum is $\mathbb A^1$. Let $m$ be its normal
meridian and let $n$ be the adjacent branch meridian at that node.
A punctured normal bundle over $\mathbb A^1$ has fundamental
group $\Z\langle m\rangle$. Hence the two-dimensional node
lattice maps into that cyclic group before mapping to
$\pi _1(U)$, and has a nonzero kernel. Thus neither type of
vertex can occur. A graph not equal to its core has a
genus-zero vertex of valency at most one at the first deletion,
so this proves the claim.
\end{proof}

\begin{theorem}[Common corner refinement of core boundaries;
\status{Proved}]\label{P3a-core}
Two marked normal-crossing surface completions in the same
meromorphic class, each with boundary graph equal to its core,
have a common refinement obtained from both by corner blow-ups.
In particular, every good completion in one meromorphic class
is connected to every other by corner blow-ups and blow-downs
through good pairs.
\end{theorem}
\begin{proof}
Resolve the marked bimeromorphic map by point blow-ups supported
on the boundary, obtaining $W\to X_1,X_2$ with normal-crossing
total boundary. Resolution of a meromorphic map of smooth
surfaces and factorization of a proper birational surface
morphism into point blow-ups apply in the analytic as well as
the algebraic category; for the algebraic factorization see
\citet[Tag 0C5Q, Lemma 54.17.1]{StacksSurfaceFactorization2026}.
A point blow-up at a node subdivides an edge; at a smooth
boundary point it attaches a rational leaf. Consequently the
strict transform of every vertex of either endpoint graph lies
in the core of the graph of $W$. Every vertex outside that core
is exceptional for both projections.

If there is such a vertex, choose the one created last in a
point-blow-up factorization of $W\to X_1$. A later blow-up touching
it would create a later noncore vertex, whether at a smooth point
or at a node. Thus it has square $-1$ and valency one or two.
Contract it. The surface remains smooth and the boundary remains
normal crossing; valency two may create an allowed self-node.
Both projections descend, since they contract that curve. Repeat.
This terminates because each contraction removes one vertex.

The resulting common refinement has no fringe. In a point-blow-up
factorization from either endpoint, any blow-up at a smooth
boundary point would create a fringe which later point blow-ups
could only subdivide or enlarge, never turn into core. Hence
both factorizations consist exclusively of corner blow-ups.
For good endpoints, Theorem \ref{P2b} proves that every pair in
these two factorizations is good.
\end{proof}

\begin{theorem}[A nonemptiness test from any completion in one class;
\status{Proved}]\label{P3a-existence}
Fix any marked normal-crossing completion $(Y,B)$ of a noncompact
smooth complex surface $U$. There is a good completion in its
meromorphic class if and only if all the following tests hold.
First, $U$ is aspherical. Second, every connected component of
$B$ has nonempty core, and its entire fringe can be contracted
to smooth points, by successively contracting rational $(-1)$
tips or bivalent curves outside the core while retaining normal
crossings. Third, on the resulting core pair $(Y_c,B_c)$ the
meridian of every boundary component has infinite order in
$\pi _1(U)$ and the rank-two meridian map at every node is
injective. The contraction of the entire fringe, when possible,
is unique over its complement. Every good completion in this
class is precisely a corner modification of $(Y_c,B_c)$.
\end{theorem}
\begin{proof}
If the tests hold, all boundary curve strata of $(Y_c,B_c)$ are
aspherical: a positive-genus curve remains aspherical after
punctures are removed, and every rational curve has at least
two normalization punctures. Point strata are contractible.
Together with the first and third tests this is exactly
goodness. The common-refinement theorem then gives every good
completion in the stated class.

For necessity, resolve a bimeromorphic map from $Y$ to a good
completion $X$. On the resulting $W$, each fringe tree is
exceptional over $X$, because $X$ is core. Starting with the
last created noncore curve in a factorization over $X$ contracts
the whole fringe to smooth points and leaves a corner
refinement of $X$. Consider now a point-blow-down factorization
$W\to Y$. A core curve carrying an attached fringe cannot be
contracted while normal crossings are retained. A positive-genus
curve cannot be an exceptional $(-1)$ curve, and a rational core
vertex has core valency at least two, so an attached fringe
raises its valency to at least three. All core contractions in
this factorization are therefore disjoint from the remaining
fringe, and commute with its contraction. A noncore contraction
merely contracts part of an already smoothly contractible
exceptional tree. Induction through the factorization proves
that the fringe already visible on $Y$ is smoothly contractible.
It also shows that its contracted core model lies in the corner
class of $X$. Blow-ups preserve whether a graph component has
empty core, so no boundary component can have empty core.

Theorem \ref{P2b} transfers goodness from $X$ to this core model,
giving the first and third tests. Finally, a proper contraction
of a fixed exceptional forest to smooth points is unique: its
target structure sheaf is the direct image of the source
structure sheaf, and the resulting map is an isomorphism off
the forest. The two targets therefore identify across the
contracted points. This is an exact test from an arbitrary
starting completion, not a claim that asphericity or group
injectivity is algorithmically decidable.
\end{proof}

A \emph{minimal good completion} here means that no nonidentity
boundary contraction to another smooth normal-crossing good
completion is possible. This is boundary minimality; it is not
minimality of the underlying compact surface.

\begin{theorem}[Existence and the precise uniqueness criterion;
\status{Proved}]\label{P3a-minimal}
Every nonempty meromorphic class of good surface completions
contains a minimal good completion. Such a completion is minimal
exactly when its boundary has no smooth rational bivalent
$(-1)$ component. Let $(X,D)$ be one minimal good completion
in the class. It is the unique minimal marked completion in
that class if and only if
\[
 \begin{cases}
 C^2\leq-2 &\text{for every smooth rational bivalent component }C,\\
 C^2\leq 2 &\text{for every singly self-nodal rational bivalent
                    component }C.
 \end{cases}
\]
The second type is an isolated rational component with one
ordinary self-node; its two normalization branches count as
its valency two. Positive-genus components and rational
components with at least three punctures impose no additional
condition.
\end{theorem}
\begin{proof}
A boundary curve contractible to a smooth point is a smooth
rational $(-1)$ curve. In a core boundary it has valency at least
two, and retaining normal crossings forces valency exactly two.
Its contraction is inverse to a corner blow-up and preserves
goodness by Theorem \ref{P2b}. Repeatedly contracting such curves
terminates because the number of components decreases. This
proves both existence and the characterization of minimality.

Let $X'$ be another minimal completion and take a common corner
refinement $W$. If the marked models differ, then in a
corner-contraction sequence $W\to X'$ at least one strict
transform of an original component of $X$ disappears. Consider
the first one to disappear. Every preceding contracted curve
was exceptional over $X$, so each preceding model still maps
to $X$. A smooth original rational bivalent component therefore
has self-intersection at most its original value until it
disappears. If that value was at most $-2$, it cannot become a
$(-1)$ curve. Positive genus and valency at least three prevent
contraction for the reasons already given.

An original self-nodal rational component must first have its
node blown up before it can be contracted as a smooth curve.
Blowing up that node decreases its self-intersection by four,
since its multiplicity there is two. While the node remains
resolved and no original component has disappeared, its square
is at most $C^2-4$. If $C^2\leq2$, this is at most $-2$, again
preventing its contraction. If the resolving exceptional is
first contracted, the component becomes singular again and
still cannot be the next smooth exceptional curve. Under the
two inequalities no first original component can disappear.
Thus every contraction from $W$ to $X'$ contracts only curves
exceptional over $X$; since both endpoints are minimal, the
marked models agree.

Conversely, a smooth rational bivalent component in a minimal
pair with square $a\geq0$ can be made $(-1)$ by $a+1$ successive
corner blow-ups adjacent to it. Contract that component, then
contract any remaining bivalent $(-1)$ components until a
minimal pair is reached. Its divisorial valuation has
disappeared and cannot reappear under further contractions,
so the resulting marked minimal pair differs from $X$.
For a self-nodal rational component with square $a\geq3$,
first blow up its node. Its square is $a-4\geq-1$; perform
further adjacent corner blow-ups if necessary, then the same
contraction argument applies. This proves necessity of both
inequalities and their sufficiency.
\end{proof}

\begin{corollary}[The exact analytic scope; \status{Proved}]
\label{P3a-analytic-scope}
For an arbitrary smooth complex surface $U$, its set $\mathcal C(U)$
of good marked compactifications is the disjoint union, over
marked meromorphic completion classes admitting a normal-crossing
model, of the corner classes passing Theorem \ref{P3a-existence}.
It is empty if there is no such class. It has a unique minimal
member if and only if exactly one class passes and the inequalities
of Theorem \ref{P3a-minimal} hold there. For a fixed algebraic $U$
and algebraic markings there is only one birational completion
class, so the preceding criterion is global without this union.
If $U$ is compact, the only completion is $(U,\varnothing)$,
and it is good exactly when $U$ is aspherical.
\end{corollary}
\begin{proof}
Every boundary blow-up or blow-down is bimeromorphic and respects
the marking, so it stays in one meromorphic class. Theorem
\ref{P3a-core} and the nonemptiness and uniqueness tests give
exactly the assertions within each such class. In the algebraic
case the identity of the dense open $U$ always gives a birational
map, so there is one class. A compact open subset of a connected
Hausdorff compact surface is closed as well, and hence equals
that surface, proving the last assertion.
\end{proof}

\subsection{The six surfaces and the analytic qualifications}

\begin{lemma}[Extension of maps to a hyperbolic affine curve;
\status{Proved}]\label{P3a-Picard}
Let $D$ be a normal-crossing divisor in a complex manifold $X$,
and let $C$ be a smooth affine complex algebraic curve with
$2g(\overline C)-2+\#(\overline C\setminus C)>0$. Every holomorphic
map $X\setminus D\to C$ extends holomorphically to $X\to\overline C$.
\end{lemma}
\begin{proof}
First choose a finite etale cover $C'\to C$ whose smooth compactification
has genus at least two. If $g(\overline C)\geq2$, take the identity.
In genus zero choose three punctures as $0,1,\infty$ and use
$y^5=x(x-1)$, restricting away from all the other punctures. Its
degree-five compactification has ramification index five at those
three points, and Riemann--Hurwitz gives genus two. In genus one,
use the degree-three permutation cover with handle monodromies
$(12),(23)$, monodromy the inverse of their commutator at one
puncture, and trivial monodromies at all other punctures. The
relation in the punctured torus group is satisfied, and the action
is transitive. Compactification has one ramification point of
index three; Riemann--Hurwitz again gives genus two. Existence of
the corresponding holomorphic cover is the covering-space form
of the Riemann existence theorem.

Work locally on $(\Delta^*)^r\times\Delta^k$ near $D$. The pullback
of this finite cover has finite monodromy, so a finite coordinate
Kummer cover kills it and permits a lift of the map to $C'$.
Compose with $C'\subset\overline C'$. Give $\overline C'$ its
compact hyperbolic metric, with positive systole. For each
punctured coordinate, hold the other coordinates fixed and take
a coordinate circle of radius tending to zero. Its length in
the complete Poincare metric of $\Delta^*$ tends to zero.
Schwarz--Pick makes its image shorter than the systole for a
sufficiently small radius. Thus that generator of
$\pi _1((\Delta^*)^r\times\Delta^k)=\Z^r$ maps trivially to
$\pi _1(\overline C')$. This works for every coordinate generator.

The lifted map therefore lifts again to the universal cover
$\Delta\to\overline C'$. This holomorphic disc-valued function is
bounded, so successive applications of the removable-singularity
theorem extend it over the coordinate divisor to $\Delta^{r+k}$.
The maximum-modulus principle keeps its image in the open disc:
an attained boundary value would force constancy with modulus one,
contrary to its values on the original open set. Composing with
$\Delta\to\overline C'\to\overline C$ gives an extension on the
Kummer cover. It is invariant under the finite deck group, since
it agrees with the original pulled-back map on a dense open set;
hence it descends holomorphically. Uniqueness on overlapping
charts glues the extensions. This direct proof also supplies the
normal-crossing version of the familiar three-punctured-sphere
extension used in \citet[p.~213]{KiernanKobayashi1973}.
\end{proof}

\begin{theorem}[The six requested computations; \status{Proved}]
\label{P3a-examples}
The following describes all good marked completions, with normal
crossings allowed to include self-nodes.
\begin{enumerate}
\item For $\C^2$, for $\PP^2$ minus a smooth conic, and for
$\PP^2$ minus an irreducible nodal cubic, $\mathcal C(U)$ is empty,
in the analytic as well as the algebraic category.

\item For $U=(\C^*)^2$ with algebraic marking, the good completions
are exactly the smooth complete toric pairs. Their minimal members
are $\PP^2$ with its coordinate triangle and the Hirzebruch surfaces
$\mathbb F_n$ with toric boundary, where $n=0$ or $n\geq2$.
In the analytic category the full list consists of those toric
pairs with arbitrary holomorphic markings, and the one-curve Hopf
pairs in Ueda's classification with arbitrary holomorphic markings,
in both cases modulo marked pair isomorphism. Every one-curve Hopf
pair is already minimal and has no other good completion in its
marked meromorphic class.

\item Put $B=\PP^1\setminus\{0,1,\infty\}$ and fix the coordinates
$(x,b)$ on $U=\C^*\times B$. The minimal analytic marked completions
are exactly
\[
 X=\PP(\mathcal O\oplus L)\longrightarrow\PP^1,\qquad
 D=S_0+S_\infty+F_0+F_1+F_\infty,
\]
where $L$ is a holomorphic line bundle on $\PP^1$ equipped with a
holomorphic trivialization $\tau$ on $B$, and the marking is
$(x,b)\mapsto x\tau(b)$ in the punctured line bundle. The classes
are precisely the relative Picard group
$\operatorname{Pic}_{\rm an}(\PP^1,B)$, with
\[
 0\longrightarrow \mathcal O^*(B)/\C^*
 \longrightarrow\operatorname{Pic}_{\rm an}(\PP^1,B)
 \longrightarrow\Z\longrightarrow0.
\]
The algebraic version is
$\operatorname{Pic}_{\rm alg}(\PP^1,B)\cong
\operatorname{Div}_{\{0,1,\infty\}}(\PP^1)\cong\Z^3$.
Every good completion is a corner refinement of one of these
minimal pairs.

\item If $U=C_1\times C_2$ for two smooth affine hyperbolic curves,
every analytic marked completion belongs to the meromorphic class
of $\overline C_1\times\overline C_2$. The unique minimal good
marked completion is the product completion with its product
boundary, and every other good completion is obtained from it by
corner blow-ups. The same statement holds algebraically.
\end{enumerate}
\end{theorem}
\begin{proof}
For any noncompact surface a compactification has nonempty
boundary. Its generic normal lattice is $\Z$. Thus finite
$\pi _1(U)$ is an absolute obstruction, irrespective of the
analytic or algebraic category. The three groups in the first
item are respectively $1$, $\Z/2$, and $\Z/3$.

Here are direct calculations for the latter two. The identification
$\operatorname{Sym}^2(\PP^1)\cong\PP^2$ identifies the diagonal
with a smooth conic. The ordered configuration space of two
distinct points of $S^2$ projects to $S^2$ with contractible fibre
$S^2\setminus\{\mathrm{point}\}\cong\C$, and is simply connected.
The free interchange involution gives the unordered complement,
so its group is $\Z/2$. For a nodal cubic choose its affine form
$y^2=x^2(x+1)$. Projection to $x$ has two generic punctures in
the fibre. At $x=-1$ its braid monodromy is a half twist and
identifies their meridians; at $x=0$ it is a full twist and
imposes commutation. Van Kampen therefore gives $\Z$ for the
affine complement. Adding the line at infinity kills its
meridian, which is the product of three generic cubic meridians,
hence the third power of the generator. There is only this
additional normal relation, because the added divisor is
connected and smooth in the complement. The projective group
is $\Z/3$.

For the algebraic torus, the standard toric completion is good:
all strata are tori, and a smooth cone's primitive rays extend
to a lattice basis, so its normal meridians form an injective
coordinate sublattice of $\pi _1(U)=\Z^2$. Every algebraic marked
completion belongs to its birational class. Theorem
\ref{P3a-core} gives a common corner refinement with any good one.
Starting from the standard toric pair, these blow-ups are toric,
and their corner contractions descend the torus action.
Consequently the other pair is toric too. The minimal smooth
complete toric surface fans are the three-ray projective-plane
fan and the four-ray Hirzebruch fans, with the $(-1)$ boundary
of $\mathbb F_1$ contractible. This also follows from the explicit
cycle reduction in \citet[Proposition~3, pp.~89--90]{Ueda1979}.

For completeness we address all analytic torus completions,
rather than silently assuming they are bimeromorphic to the
standard algebraic marking. Ueda's Theorem~2 lists rational
surfaces, one-curve Hopf surfaces and their boundary blow-ups,
and the Serre ruled surface over an elliptic curve and its
boundary blow-ups \citep[p.~81]{Ueda1979}.
In the rational case $D$ is connected, since $(\C^*)^2$ has one
end. Poincare--Lefschetz duality and the homology sequence of
$(X,D)$ give a surjection
\[
 H^2(U,\Q)=H_2(X,D;\Q)\longrightarrow H_1(D,\Q),
\]
because $H_1(X,\Q)=0$. The left group has dimension one. Since
\[
 b_1(D)=2\sum_i g(\widetilde D_i)+b_1(\text{incidence graph}),
\]
all components are rational. A connected core graph with first
Betti number at most one is a cycle. Resolve a self-node and,
if needed, subdivide a two-edge cycle by corner blow-ups.
The resulting simple cycle satisfies Ueda's Proposition~3 and
reduces by corner contractions to the projective-plane or
Hirzebruch toric pair. The original pair is therefore toric
as well; its given holomorphic marking need not be algebraic.

A Hopf pair on Ueda's list has one smooth elliptic boundary.
On its primary Hopf cover the complement is $\C\times\C^*$,
and a small normal meridian is the generator of its $\Z$.
That cover is a genuine covering of the complements, so this
meridian has infinite order downstairs. The same argument
passes to secondary Hopf quotients. The open torus and the
elliptic boundary are aspherical, proving goodness. Any
boundary blow-up adds a rational fringe and is bad. Since
there are no nodes, the common-corner theorem also proves
uniqueness in its meromorphic class.

For the Serre ruled surface, the complement is an affine-line
bundle over an elliptic curve. Both that bundle and its
$\PP^1$ compactification induce the same fundamental group
as the base. Hence the normal boundary meridian dies. Boundary
blow-ups do not change the generic meridian of its strict
transform, so none can be good. This proves the complete
analytic torus list.

We next treat $\C^*\times B$. By Lemma \ref{P3a-Picard}, the
projection to $B$ extends to a holomorphic map
$p:X\to\PP^1$ on every analytic completion. Its general fibre
is the compactification of $\C^*$, hence $\PP^1$, and has
exactly two horizontal boundary points. The given coordinate
$x$ labels its two ends consistently on $B$, so these horizontal
branches have trivial permutation monodromy. Their closures
are two sections: on normalization each maps to the base
with degree one. In particular each is smooth. The Stein
factorization has degree one because the general fibre is
connected, so all fibres of $p$ are connected.

This already proves rationality without assuming a minimal
underlying surface. For one of the sections $S$, the coherent
sheaf $p_*\mathcal O_X(S)$ is torsion free on $\PP^1$, hence
locally free, and has generic rank two. Its evaluation map
defines a meromorphic map from $X$ to the projectivization of
this rank-two bundle, restricting to the degree-one complete
linear system on the general fibre. It is therefore bimeromorphic.
A holomorphic bundle over $\PP^1$ splits, so that projectivization
is a Hirzebruch surface. Resolving this bimeromorphic map and
using point factorization shows that $X$ is a smooth rational
projective surface. These auxiliary contractions are used to
prove rationality; they are not asserted to preserve the
original open $U$. This avoids an invalid appeal to a theorem
assuming the underlying compact surface was already minimal.

Consequently $p$ and $D$ are algebraic in this projective model.
Choose a rational fibre coordinate $t$ having its generic zero
and pole at the two sections. Any zero or pole of $t$ on
$U$ is vertical. The fibres of $U\to B$ are smooth irreducible
copies of $\C^*$, so
$\operatorname{div}(t)|_U$ is a finite sum of entire fibres.
Since $\operatorname{Pic}_{\rm alg}(B)=0$, multiply $t$ by a
rational base function to remove all these vertical divisors.
The resulting $t$ is an algebraic unit on $U$. On every fibre,
write it in the original coordinate as
$t=a(b)x^{k_b}$. This follows because each algebraic fibre has
compactification $\PP^1$ with two points deleted. Its winding
number $k_b$ is locally constant in $b$, and generically is
$1$ or $-1$. Thus, after inversion if needed, $k_b=1$
everywhere and $a$ is a nowhere-zero holomorphic function on
$B$. The map $(p,t):U\to B\times\C^*$ is an algebraic etale
bijection, hence an algebraic isomorphism. This supplies an
auxiliary algebraic marking; the given analytic marking
differs from it by the unit $a(b)$.

In this auxiliary algebraic marking the core theorem supplies
a common corner refinement with the product completion.
All corner operations lie over $0,1,\infty$. Neither horizontal
section can be contracted because $p$ is nonconstant on it.
In a reducible fibre arising from corner operations, the
rational chain has a bivalent $(-1)$ component until it is
irreducible; this is also the usual point-blow-up description
of a fibre of a ruled surface. Thus a minimal good model has
only smooth irreducible fibres and is a $\PP^1$ bundle with
two disjoint sections. These sections split its rank-two
bundle, giving $\PP(\mathcal O\oplus L)$. Its sections have
valency three and its three boundary fibres have square zero,
so it is boundary-minimal. Conversely each such pair is good:
over $B$ it is the marked product, the section meridians are
the nonzero $\Z$ factor, and the fibre meridians project to
nontrivial peripheral elements of $\pi _1(B)$; at their
intersections this projection makes the two meridians independent.

A marked isomorphism must preserve $p$, the zero end, and the
infinity end. It is exactly a line-bundle isomorphism carrying
$\tau$ to $\tau'$. Hence its classes are relative Picard classes.
The exact sequence displayed in the statement follows by
restricting units and line bundles from $\PP^1$ to $B$;
$\operatorname{Pic}_{\rm an}(B)=0$ because $B$ is a Stein
Riemann surface with $H^2(B,\Z)=0$. In the algebraic case a
trivialization extends to a rational section, whose divisor
is supported at the three punctures, yielding $\Z^3$.
There is no inversion quotient with the fixed marking in
the statement: $x$ labels $0$ and $\infty$. Inversion becomes
an equivalence only if one allows the marking itself to
change by $x\mapsto x^{-1}$. Holomorphic units such as
$\exp(h(b))$ explain why the analytic classification is
strictly larger than the algebraic $\Z^3$ list.

Finally, for $C_1\times C_2$, both projections extend by the
Picard lemma. Their product is a proper holomorphic map
$X\to\overline C_1\times\overline C_2$ restricting to the
specified identity on $U$, so it is bimeromorphic. The product
pair is good: curve strata are the hyperbolic affine factors,
and normal meridians are peripheral elements in independent
free-group factors. A rational boundary component has at
least three punctures; all other components have positive
genus. Thus there are no smooth or self-nodal rational
bivalent components. Theorem \ref{P3a-minimal} proves uniqueness
of the minimal model, and Theorem \ref{P3a-core} gives all its
corner refinements.
\end{proof}

\begin{proposition}[Failures of the unrestricted analytic and stack
comparison; \status{Refuted}]\label{P3a-scope}
The assertion that every two good compactifications of a smooth
complex surface are connected by stratal blow-ups and blow-downs
is false if arbitrary analytic markings are allowed. It is also
false for smooth Deligne--Mumford compactifications if the only
allowed operations are ordinary stratal blow-ups and blow-downs.
Even existence in a specified algebraic compactification class
need not agree with existence for the underlying analytic surface.
\end{proposition}
\begin{proof}
Fix $0<|\alpha|<1$, and let
\[
 H=(\C^2\setminus\{0\})/
      \langle(z,w)\mapsto(\alpha z+w,\alpha w)\rangle,\qquad
 E=\{w=0\}/\langle g\rangle.
\]
On $w\ne0$, put $u=\alpha z/w$ and $v=w$. Then
$g(u,v)=(u+1,\alpha v)$. Choosing a logarithm of $\alpha$,
the invariant coordinates
\[
 A=e^{2\pi i u},\qquad B=v\,e^{-u\log\alpha}
\]
identify $(H\setminus E)$ biholomorphically with $(\C^*)^2$.
The pair is good by the Hopf argument above. It has $b_1(H)=1$,
whereas a toric rational completion has $b_1=0$.
Blowing up points preserves $b_1$, so these two good analytic
pairs cannot be connected by the requested operations.
Their dual complexes are respectively a point and a circle.
Even on a fixed toric surface, changing the marking by
$(x,y)\mapsto(x,e^x y)$ need not give the same meromorphic
class: an alleged meromorphic extension would restrict on a
generic $y$-section to the meromorphic continuation of $e^x$
across $x=\infty$, which is impossible.

For the stack assertion, start with a good toric pair and
take a nontrivial root of a boundary divisor. Its strata are
finite gerbes over torus strata and are aspherical in the
stack sense. On an atlas its meridian is a nonzero multiple
of the old primitive meridian, so the pair is good.
Triviality of all inertia is preserved by an ordinary
representable blow-up in both directions. One direction follows
because blowing up an algebraic space again has trivial inertia.
For the converse, if a nonidentity finite stabilizer element $g$
fixes a point outside the centre, it persists upstairs. If the
point lies in the centre, the blow-up fibre is the projective
space of its normal representation. Over $\C$, $g$ has an
eigenline, so some point of that fibre still has nontrivial
stabilizer. Thus no zigzag of these blow-ups can connect the
root pair to the toric pair with trivial inertia. Generic
inertia on a particular boundary divisor is not claimed to be
invariant: exceptional divisors can acquire such inertia.

For the last assertion let
$H_q=(\C^2\setminus\{0\})/\langle(z,w)\mapsto(qz,qw)\rangle$,
$0<|q|<1$, and delete $E_1=\{w=0\}/q^\Z$.
Then
\[
 H_q\setminus E_1\cong E_q\times\C,\qquad
 E_q=\C^*/q^\Z,\qquad
 (z,w)\longmapsto([w],z/w).
\]
The boundary $E_1$ is elliptic. Near it, use $[z]$ and
$\xi=w/z$; its normal meridian maps to
$([\xi z],1/\xi)$. The angular loop in $E_q$ is nonzero
and of infinite order, so this is a good analytic completion.
For the fixed algebraic surface $E_q\times\mathbb A^1$, no
algebraic good completion exists. In any proper algebraic
completion, the map $\mathbb A^1\to U$ given by one affine
fibre extends to $\PP^1\to X$. The sole missing point has
positive contact order with at least one boundary branch.
A sufficiently large circle bounds the finite affine disc
inside $U$, but maps to a nonzero full normal-meridian
vector at that missing point. This contradicts goodness.
Equivalently the standard algebraic completion has a
boundary meridian which dies, and the core test forbids all
other algebraic completions. The Hopf and algebraic ruled
surfaces are not bimeromorphic: their first Betti numbers
are one and two. Thus a meromorphic-class qualification
is indispensable even for nonemptiness.
\end{proof}

\subsection{Elementary fibrations, stable boundaries, and arithmetic inertia}

An elementary hyperbolic fibration here means a smooth family of affine
curves obtained from a smooth proper family by removing a nonempty
finite etale marked divisor, with $2g-2+n>0$. Analytically it is a
locally trivial bundle with fibre a punctured compact surface; its
fundamental group is a finitely generated nonabelian free group.
An iterated tower starts from a point. Base points and paths are
specified when homomorphisms, rather than their conjugacy classes,
are wanted.

\begin{theorem}[Horizontal inertia over an arbitrary base;
\status{Proved}]\label{P3b-discrete}
For an elementary hyperbolic fibration $F\to U\to B$ with connected
base, the fibre map $\pi _1(F)\to\pi _1(U)$ is injective. Every
horizontal puncture meridian has infinite order in $\pi _1(U)$.
Neither assertion requires $B$ to be aspherical. An iterated
hyperbolic tower is aspherical.
\end{theorem}
\begin{proof}
The possible obstruction to fibre injectivity is the connecting map
$\partial:\pi _2(B)\to\pi _1(F)$. Pull the bundle back along a
representative sphere. Its clutching loop is a loop $h_t$ of
self-homeomorphisms of $F$, based at the identity. At a chosen point
$x$, the loop $h_t(x)$ represents $\partial$. For any based loop
$\gamma(s)$ in $F$, the square $h_t(\gamma(s))$ supplies a homotopy
showing that $h_t(x)$ commutes with $\gamma$. Hence the image of
$\partial$ is central. The centre of a nonabelian free group is
trivial: a reduced word commuting with a free generator is its
power, and the cyclic subgroups of two different generators
intersect trivially. Thus $\partial=0$.

A puncture circle in a hyperbolic affine curve is nontrivial.
For at least two punctures it survives in homology; for one
puncture and positive genus it is the nontrivial reduced product
of handle commutators in the free group. A free group is
torsion-free, so every such meridian is infinite. The fibre
injection proves the horizontal assertion. Finally, the homotopy
exact sequence of a bundle with aspherical base and fibre gives
asphericity of its total space. Induction proves the tower
assertion.
\end{proof}

\begin{lemma}[Free normal subgroups and profinite completion;
\status{Proved}]\label{P3b-completion}
If $1\to N\to G\to Q\to1$ is an extension and $N$ is a finitely
generated nonabelian free group, then
\[
 1\longrightarrow\widehat N\longrightarrow\widehat G
 \longrightarrow\widehat Q\longrightarrow1
\]
is exact. Every nontrivial cyclic subgroup of $N$ induces an
embedding $\widehat\Z\to\widehat N$.
\end{lemma}
\begin{proof}
We first prove $Z(\widehat F_r)=1$ for $r\ge2$, rather than
assuming that discrete centre-freeness survives completion.
Let $F_r\to A$ be any finite quotient with kernel $K$, and choose
a prime $p\nmid |A|$. The regular covering of a rose with $r$
edges has cellular exact sequence
\[
 0\longrightarrow H_1(K;\mathbf F_p)
 \longrightarrow\mathbf F_p[A]^r
 \longrightarrow\mathbf F_p[A]
 \longrightarrow\mathbf F_p\longrightarrow0.
\]
Averaging over $A$ splits every sequence of these modules.
Consequently
\[
 H_1(K;\mathbf F_p)\cong
 \mathbf F_p\oplus\mathbf F_p[A]^{r-1},
\]
on which $A$ acts faithfully. The finite group
$F_r/[K,K]K^p$ is an extension of $A$ by this module, with this
conjugation action. The image of a central element of
$\widehat F_r$ in that finite group acts trivially on its kernel;
its image in $A$ is therefore $1$. Since $A$ was arbitrary,
the central element is $1$.

For a topologically finitely generated profinite group, its
continuous automorphism group is profinite. Indeed, intersections
of all open subgroups of index at most $n$ are characteristic
open subgroups forming a basis, and compatible automorphisms of
the resulting finite quotients form a closed inverse-limit group.
Conjugation gives
\[
 G\longrightarrow\operatorname{Aut}(N)
 \longrightarrow\operatorname{Aut}(\widehat N),
\]
which extends to $\widehat G$. Its composite with
$\widehat N\to\widehat G$ is the inner-automorphism map of
$\widehat N$. The centre-freeness just proved makes this
composite injective, proving $\widehat N\hookrightarrow\widehat G$.
The image is the closure of $N$, and quotienting $\widehat G$
by that closure gives $\widehat Q$, since both classify precisely
the finite quotients of $G$ killing $N$. This proves exactness
without a goodness hypothesis on $Q$.

For the cyclic assertion, represent a nontrivial conjugate of
$w\in F_r$ by a cyclically reduced word and its immersed labelled
circle in the rose. Complete that finite immersion to a finite
cover: for each edge label, pair the finitely many missing
incoming and outgoing edges. Their numbers agree. The original
circle is then an embedded cycle in a finite covering graph.
Its homology class is primitive, so the covering subgroup $H$
admits a homomorphism $H\to\Z$ sending $w$ to $1$.
For any $n$, take the kernel of $H\to\Z/n$ and then its core
in $F_r$. The resulting finite quotient of $F_r$ gives $w$
order divisible by $n$. Thus the topology induced on
$\langle w\rangle$ contains every congruence topology $n\Z$,
which is exactly the injectivity of
$\widehat\Z\to\widehat F_r$.
\end{proof}

\begin{corollary}[Horizontal profinite injection and towers;
\status{Proved}]\label{P3b-horizontal-profinite}
In Theorem~\ref{P3b-discrete}, every horizontal meridian induces
an injection $\widehat\Z\hookrightarrow\widehat{\pi _1(U)}$,
over any connected base. The group of an iterated elementary
hyperbolic tower is residually finite.
\end{corollary}
\begin{proof}
The homotopy sequence and the preceding theorem give a group
extension with kernel $\pi _1(F)$, so apply
Lemma~\ref{P3b-completion} twice, first to the fibre and then to
its nontrivial peripheral cyclic subgroup. For residual
finiteness, induct on the tower: an element nontrivial in the
base is separated by a finite quotient of the base, and a
nontrivial element in the fibre survives in its embedded
profinite completion. The latter group is residually finite
by the same finite-cover construction. This proves the induction.
The related goodness and completion statements for elementary
anabelian towers appear in
\citet[Lemmas 10.9 and 10.11]{Bresciani2021}; here the affine
free-fibre argument is supplied explicitly, and the horizontal
statement has the stronger arbitrary-base scope.
\end{proof}

\begin{theorem}[The corrected vertical statement;
\status{Proved}]\label{P3b-vertical}
Suppose the fibration extends to a compatible stable pointed
family over an NC base completion, after the indicated finite
root or quotient base changes, and take a smooth toroidal model
of the total boundary. Assume:
\begin{enumerate}
\item the base's open strata are aspherical;
\item its discrete full meridian lattices inject in $\pi _1(B)$;
\item for the profinite assertion, its completed full meridian
lattices also inject in $\widehat{\pi _1(B)}$.
\end{enumerate}
Then the total completion is good, and under the third hypothesis
its full completed meridian lattices inject as well. These
hypotheses hold inductively for recursively compatible
stable/toroidal completions of elementary hyperbolic towers
whose initial curve completion has its usual peripheral
injections. Stable extension alone does not imply either
vertical conclusion for an arbitrary chosen base completion.
\end{theorem}
\begin{proof}
Stable curve components, punctured at their nodes and markings,
are aspherical. A rational component has at least three special
points; an elliptic component has a special point unless it is
a whole unpointed genus-one curve, excluded from the hyperbolic
stable pointed problem. Toroidal resolution can additionally
insert rational bridges with two punctures, which are
$\C^*$. Compatibility includes constancy of the fibre incidence
type on each base stratum. The total open strata are therefore
bundles of these curves, points, or tori over finite incidence
covers of base strata. Their homotopy sequences prove their
asphericity under the first hypothesis; the open total space is
aspherical when the open base is, by Theorem~\ref{P3b-discrete}.

Consider a local total normal lattice $A$ and the corresponding
base lattice $A_B$. Projection induces an integral map $M:A\to A_B$.
At a smooth vertical point it has zero kernel; at a horizontal
marking its kernel is the fibre's puncture circle. At a node
the basic local equation is $xy=t$, for which
\[
 \Z^2\longrightarrow\Z,\qquad (u,v)\longmapsto u+v,
\]
and the kernel generator $(1,-1)$ is the vanishing cycle.
After finite base change the equations are monomial; smooth
toroidal charts replace the displayed lattices by primitive
sub-lattices and finite-index lattices. Thus at every point
the kernel has rank at most one and, if present, maps to a
nonzero power of a puncture or vanishing cycle in $\pi _1(F)$.

The relevant vanishing cycle is nontrivial. A smoothing neck
could bound a disc with no puncture only if the special fibre
had an unstable unmarked rational tail on that side; stable
reduction removes it. A peripheral neck enclosing one puncture
is still nontrivial, as are the necks in inserted two-punctured
bridges. All other necks are essential surface curves.
This verifies the kernel claim also in the separating case.

Now use the commutative diagram with exact rows
\[
\begin{array}{ccccc}
 \ker M&\longrightarrow&A&\longrightarrow&A_B\\
 \downarrow&&\downarrow&&\downarrow\\
 \pi _1(F)&\longrightarrow&\pi _1(U)&\longrightarrow&\pi _1(B).
\end{array}
\]
The right vertical arrow is injective by hypothesis, and the
left arrow is injective because a nontrivial free-group element
has infinite order. Any element killed by the middle arrow
first lies in $\ker M$ and then is zero. This proves the full
discrete lattice assertion.

Complete the bottom row using Lemma~\ref{P3b-completion}.
The kernel of the completed integral matrix $M$ is the
completion of its integral kernel: the image is a free abelian
group, the short exact sequence with that image splits, and
its inclusion in $A_B$ is still injective after completion.
The third hypothesis and the cyclic part of the same lemma
now give the identical diagram chase with $\widehat A$.
Starting from a punctured hyperbolic curve, these discrete and
completed peripheral injections hold by the finite-cover
argument above. Roots restrict to finite-index lattices, and
stratal subdivisions use unimodular charts, so the induction
survives all the specified compatible operations. The theorem
is conditional on such a compatible stable/toroidal completion;
it does not silently assert that every arbitrarily chosen
completion of a tower is compatible.
\end{proof}

\begin{example}[Stable extension on a bad base; \status{Refuted}]
\label{P3b-vertical-counterexample}
Take $B=(\PP^1\setminus\{0,1,\infty\})^2$ with its product
completion, and blow up a smooth point of one boundary component.
At the new boundary node the strict-transform and exceptional
meridians both map to the same old meridian. Their difference
is killed. For any hyperbolic affine curve $F$, the constant
stable marked family $\overline F\times Y\to Y$ over this
blown-up base is already stable; its open total space is the
genuine three-stage hyperbolic tower $F\times B$.
Nevertheless the same rank-two vertical kernel persists in
$\pi _1(F)\times\pi _1(B)$, and in its profinite completion.
This refutes the unrestricted vertical assertion even within
the requested class of towers. It also illustrates why
injectivity of individual meridians is inadequate.
\end{example}

\begin{corollary}[Characteristic-zero geometric inertia;
\status{Proved}]\label{P3b-profinite}
Let the data be defined over a field $K$ finitely generated over
$\Q$, and choose $\overline K\hookrightarrow\C$. Under the
hypotheses just established, horizontal geometric inertia,
and the full vertical inertia in the compatible tower case,
inject into $\pi _1^{\mathrm{\acute et}}(U_{\overline K})$:
\[
 T^{I(x)}=\widehat\Z(1)^{I(x)}
       \hookrightarrow\pi _1^{\mathrm{\acute et}}(U_{\overline K}).
\]
These are homomorphisms up to the conjugations determined by
geometric base-point paths; no preferred trivialization of
the Tate module is asserted.

If the local punctured-neighbourhood extension
$1\to I_x\to E_x\to G_{K(x)}\to1$ maps to
$1\to\pi _1(U_{\overline K})\to\pi _1(U)\to G_K\to1$
and the chosen $G_{K(x)}\to G_K$ is injective, then
\[
 \ker(E_x\to\pi _1(U))
 =\ker(I_x\to\pi _1(U_{\overline K})).
\]
Thus in that case full local arithmetic injectivity is
equivalent to geometric inertia injectivity. This qualification
holds for a closed point after the corresponding finite field
extension, with compatible separable closures. It must not
be omitted for a transcendental residue field.
\end{corollary}
\begin{proof}
Riemann existence identifies the geometric etale fundamental
group with the profinite completion of the complex topological
one; the precise comparison is
\citet[Exposé XII, Théorème 5.1 and Corollaire 5.2]{SGA1}.
In an NC chart the compatible Kummer covers of the punctured
coordinates give $T^{I(x)}$. Under the chosen complex embedding
and a compatible system of roots of unity, this is precisely
the completed normal lattice used above. The proved profinite
injections therefore give the assertion. Changing paths
conjugates the maps; Galois acts cyclotomically and may permute
branches, which is why the intrinsic Tate-module notation is
retained.

For the final statement, project any element of
$\ker(E_x\to\pi _1(U))$ to $G_{K(x)}$. Its image in $G_K$ is
trivial, hence its image in $G_{K(x)}$ is trivial by the stated
hypothesis. It belongs to $I_x$, where the two kernels agree.
The reverse inclusion is immediate. For a transcendental
residue field the Galois-quotient map need not be injective,
so this argument gives no such equivalence. None of these
characteristic-zero arguments treats wild inertia in positive
characteristic.
\end{proof}

\subsection{Higher-dimensional comparison: established results and the remaining question}

\begin{proposition}[An additional projective obstruction;
\status{Proved}]\label{P3c-Mori}
A good smooth projective NC pair admits no nonconstant algebraic
map $\mathbb A^1\to S$ into any of its open strata, including $U$.
For a simple normal-crossing boundary it follows that
$K_X+D$ is nef.
\end{proposition}
\begin{proof}
Let $f:\mathbb A^1\to S$ be nonconstant. Properness extends it to
$\overline f:\PP^1\to X$, with image in the closure of $S$.
If the image of infinity remains in $S$, asphericity makes this
sphere nullhomotopic in $S$, hence in $X$. But the pullback of
a Kähler form of the projective $X$ has strictly positive
integral on a nonconstant holomorphic sphere, a contradiction.

Otherwise infinity maps to a deeper stratum $T$. Pull back the
normal directions of $S$ along $\mathbb A^1$. Their bundles and branch
labelling are topologically trivial because $\mathbb A^1$ is contractible.
Choose nonzero normal displacements and a continuous positive
radius that tends to zero faster than every newly vanishing
boundary coordinate near infinity. This pushes $f(\mathbb A^1)$ into
$U$. Its large circle bounds the pushed finite disc. In a local
NC chart at $\overline f(\infty)$, however, the circle has a
nonzero meridian vector: on every branch newly acquired at
infinity the winding is the nonzero order of vanishing, with
a common orientation sign. The radius choice preserves these
windings. The coordinates in the old normal directions can be
arbitrary integers, but cannot cancel the distinct new
coordinates. Thus a nonzero element of the full lattice at
$T$ dies in $\pi _1(U)$, contrary to goodness. This proves
the first assertion, also when a normalization of the stratum
closure is needed to keep track of branches.

For a reduced SNC divisor these open strata are exactly those
in the Mori-hyperbolicity hypothesis of
\citet[Theorem 1.1]{Svaldi2019}. That theorem, applied to the
smooth projective log canonical pair $(X,D)$, gives
$K_X+D$ nef. This last step uses the cone theorem; no claim
that asphericity alone implies positivity is being made.
\end{proof}

\begin{theorem}[Birational morphisms between good projective pairs;
\status{Proved}]\label{P3c-morphism}
Let $f:(Y,E)\to(X,D)$ be a proper birational morphism of smooth
projective varieties with reduced SNC boundaries, inducing the
specified identity on $U=Y\setminus E=X\setminus D$. If both pairs
are good, then $f$ is logarithmically etale and toroidal. Moreover
\[
 K_Y+E=f^*(K_X+D).
\]
This asserts toroidality of the morphism, not a factorization of
it solely by smooth stratal blow-ups.
\end{theorem}
\begin{proof}
The natural logarithmic differential map
\[
 f^*\Omega_X^1(\log D)\longrightarrow\Omega_Y^1(\log E)
\]
is defined everywhere. Indeed, in boundary coordinates each
$f^*x_i$ is a unit times a monomial in the local equations of
$E$, so its logarithmic differential has only logarithmic poles.
It is an isomorphism on $U$. Its determinant has an effective
divisor $A$, and
\[
 K_Y+E=f^*(K_X+D)+A.
\]
At the generic point of every nonexceptional prime divisor a
proper birational morphism between smooth varieties is an
isomorphism. Thus $A$ is $f$-exceptional. Proposition
\ref{P3c-Mori} makes $K_Y+E$ nef, so $A$ is $f$-nef. The
negativity lemma applied to $-A$ gives $-A\geq0$, because
$f_*(-A)=0$. With $A\geq0$ this gives $A=0$. We use the sign
convention of \citet[Lemma~3.39]{KollarMori1998}, also stated
in \citet[Lemma~1.73, author version p.~51]{Xu2025}.
The determinant has no zero, so the logarithmic differential
map is an isomorphism everywhere.

Here is the local chart implication, in addition to the abstract
criterion \citep[Proposition~3.12 and Theorem~3.5]{Kato1989}.
Let $x_1,\ldots,x_r$ be the target boundary coordinates and
$y_1,\ldots,y_s$ the source boundary coordinates at the points
in question. Write
\[
 f^*x_i=u_i\prod_{j=1}^s y_j^{a_{ij}},
 \qquad i=1,\ldots,r,
\]
with units $u_i$, and write $g_\ell=f^*x_\ell$ for the remaining
ordinary coordinates. In logarithmic bases, at the source
point the derivative has block form
$\left(\begin{smallmatrix}A&B\\0&C\end{smallmatrix}\right)$,
where $A=(a_{ij})$ has $s$ columns. Invertibility forces
$\operatorname{rank}A=s\leq r$. Choose $s$ pivot rows $I$
whose matrix $A_I$ has nonzero determinant. Locally analytically
take logarithms of their units and put
\[
 z_j=y_j\exp\bigl((A_I^{-1}\log u_I)_j\bigr).
\]
Then the pivot target coordinates are pure monomials in the
$z_j$; every other boundary coordinate is a monomial times a
unit $h_i$. Row elimination in the same logarithmic Jacobian
shows that
\[
 z_1,\ldots,z_s,\quad h_i\ (i\notin I),\quad
 g_{r+1},\ldots,g_n
\]
are ordinary local coordinates, with the $h_i$ regarded as
invertible coordinates near their nonzero values. Algebraically
one obtains the same change after adjoining finitely many
roots of units, which is etale in characteristic zero.

Consequently the chart is induced by
$\mathbb N^r\to\mathbb N^s\oplus\Z^{r-s}$, with the nonpivot
units supplying the last generators. Its group matrix has
determinant $\pm\det A_I$, hence finite kernel and cokernel
of invertible order, and the map to the associated monoid
base change is ordinarily etale in the displayed coordinates.
Kato's chart criterion gives logarithmic etaleness. These
are exactly etale-local toroidal monomial charts, establishing
toroidality rather than merely vanishing relative differentials.
\end{proof}

\begin{proposition}[Crepancy for rational comparisons;
\status{Proved}]\label{P3c-crepant}
Any two smooth projective good SNC completions of the same
algebraically marked $U$ are log-crepant: on a common smooth
NC resolution $p:W\to X$, $q:W\to Y$,
\[
 p^*(K_X+D)=q^*(K_Y+E).
\]
The effective logarithmic Jacobian divisors of the two
projections are equal and are exceptional for both.
\end{proposition}
\begin{proof}
Let $B=W\setminus U$ with reduced structure. As in the preceding
proof, the two logarithmic Jacobians give
\[
 K_W+B=p^*(K_X+D)+A_p=q^*(K_Y+E)+A_q,
\]
where $A_p,A_q$ are effective and exceptional for their respective
projections. Put
$R=p^*(K_X+D)-q^*(K_Y+E)=A_q-A_p$.
Both endpoint log canonical divisors are nef by Proposition
\ref{P3c-Mori}. Thus $-R$ is $p$-nef, and
$p_*R=p_*A_q$ is effective. The negativity lemma gives $R\geq0$.
Likewise $R$ is $q$-nef and $q_*(-R)=q_*A_p$ is effective, so
the same lemma gives $-R\geq0$. Hence $R=0$ and $A_p=A_q$.
Equality does not force this common effective divisor to vanish:
$W$ has not been assumed good. Contracting all of its positive
support compatibly with both projections is precisely a gap
which these arguments do not fill.
\end{proof}

\begin{theorem}[The complete toric comparison; \status{Proved}]
\label{P3c-toric}
Let $X_\Sigma$ and $X_{\Sigma'}$ be smooth complete toric
varieties with their full toric boundaries, marked by the
same open algebraic torus $T$. In every dimension, and in
particular for threefolds, the pairs are good and are joined
over $T$ by a finite zigzag of smooth stratal blow-ups and
blow-downs through good pairs.
\end{theorem}
\begin{proof}
Each open stratum is a torus, hence aspherical. If a smooth
cone is generated by primitive rays $v_1,\ldots,v_r$ in the
cocharacter lattice $N$, these vectors extend to an integral
basis. The local normal-meridian map is the inclusion
\[
 \Z^r\longrightarrow N=\pi _1(T),\qquad e_i\longmapsto v_i,
\]
so it is injective, both discretely and after completion.

A blow-up of the closure of the orbit associated with the
cone $\langle v_1,\ldots,v_r\rangle$ replaces it by the smooth
star subdivision with new ray $v_1+\cdots+v_r$. Thus ordinary
smooth stratal blow-ups are precisely these smooth star
subdivisions; their inverses are smooth star assemblies.
The toric weak factorization theorem supplies a finite
sequence of these subdivisions and assemblies between two
nonsingular fans with the same support. The precise theorem
used is \citet[Theorem 4.7.1]{Wlodarczyk2009}; its hypotheses
hold because both fans are nonsingular and complete in the
same $N_\R$. Every intermediate fan has that support, hence
is complete, and the preceding stratum and lattice calculation
makes every intermediate pair good. This proves the entire
requested toric case, not just preservation under a single
toric blow-up.
\end{proof}

\begin{proposition}[The weak and strong requests are different;
\status{Proved}]\label{P3c-weak-strong}
The toric zigzag theorem and the weighted strong factorization
theorem do not establish a common roof obtainable solely by
smooth stratal blow-ups.
\end{proposition}
\begin{proof}
The preceding zigzag permits blow-ups and blow-downs in
alternating order. Requiring a single common refinement
reachable from both fans solely by smooth star subdivisions
is the unweighted strong factorization problem. It is the
explicit question in
\citet[arXiv version 2, Section 7.1, Question 7.1]{AdiprasitoPak2026}.
The weighted theorem now published by those authors allows
arbitrary rational stellar subdivisions. Such a subdivision
can introduce a ray different from the sum of primitive
generators, and can pass through singular toric varieties;
it does not prove the common smooth roof assertion. The
surface common-roof theorem above therefore must not be
promoted to higher dimensions by confusing these two results.
\end{proof}

\begin{conjecture}[A restricted comparison question;
\status{Open}]\label{P3c-general}
Let $U$ be a fixed smooth algebraic threefold over $\C$, and
let $(X,D)$ and $(Y,E)$ be smooth proper good NC completions
with algebraic markings by $U$. Are they joined over $U$ by
a finite zigzag of blow-ups and blow-downs of smooth closures
of boundary strata, all intermediate pairs being good?
\end{conjecture}

The unrestricted analytic and DM-stack versions are refuted
by Proposition~\ref{P3a-scope}; even the analytic surface
version is false. The displayed scheme question is a
deliberately restricted replacement. We have proved it for
surfaces and for toric pairs. Ordinary weak factorization,
even with centres having normal crossings with the boundary,
does not answer it: Theorem~\ref{P2a} shows exactly why its
non-stratal centres are forbidden here. Propositions
\ref{P3c-Mori} and \ref{P3c-crepant} give nefness and log-crepancy,
and Theorem \ref{P3c-morphism} proves toroidality for an actual
birational morphism. They do not provide a common good resolution
of an arbitrary rational comparison. We give no proof or
counterexample for the remaining nontoric threefold question
and do not assert that its formulation constitutes a theorem
or a proof of its status in the literature.

\begin{proposition}[A failed elementary flop candidate;
\status{Proved}]\label{P3c-flop-test}
Consider an ordinary Atiyah flop between smooth projective pairs,
with exceptional curve
$C\cong\PP^1$, normal bundle $\mathcal O(-1)^{\oplus2}$.
Assume locally that exactly one boundary surface contains
$C$, that $C$ has square $-1$ on that surface, and that the
other boundary branches meet $C$ transversely at $n$ distinct
points, with no additional boundary near $C$. This construction
cannot give a good-to-good NC comparison with a non-stratal
forward blow-up.
\end{proposition}
\begin{proof}
For $n=0$, the containing open stratum contains the complete
rational curve, contrary to asphericity and Kähler positivity.
For $n=1$, its punctured part is an affine line, contrary to
Proposition~\ref{P3c-Mori}. In the common blow-up, the exceptional
quadric is $\PP^1\times\PP^1$. The containing boundary surface
meets it in a ruling fibre, whereas each transverse boundary
branch meets it in a fibre of the other ruling. After the
flop, each of those $n$ branches contains the new curve $C'$.
For $n\ge3$ this violates normal crossings along a codimension-two
curve. For $n=2$, the inverse blow-up centre $C'$ is the closure
of the intersection stratum of the two branches, hence is
stratal. If the target pair were good, its blow-up would
therefore be good by Theorem~\ref{P2b}. But the forward
blow-up has centre properly inside a single boundary surface
and is non-stratal, so is bad by Theorem~\ref{P2a}. The common
blow-up cannot be both. This excludes this precise candidate,
without treating it as a classification of all threefold
flops or a solution of Conjecture~\ref{P3c-general}.
\end{proof}
